\section{Baseline attribution and implementation notes}
\label{app:baseline-attribution}

Each of the five baselines was reimplemented inside the released evaluation
framework rather than vendored from its reference repository, for two
reasons. First, two of the five reference repositories ship no license
file: the Accenture SIFU branch and the Zoybzo FedEraser mirror. Under
default copyright law, code without an explicit license cannot be
redistributed, and the authors did not respond to clarification requests
within the paper's revision window. Second, FedRecovery has no public
reference implementation at all; the paper reports results that an
implementation produced, but the implementation itself is not available.
A reimplementation from the paper description was therefore unavoidable.

Where a reference implementation is publicly available under a permissive
license, our reimplementation was cross-checked against it to confirm that
the core update rule matches the published algorithm. Specifically,
Halimi's projected gradient ascent step and FedMemPrune's memorisation-score
mask were each verified against the authors' code on a toy problem (10-class
MNIST, five clients) before being applied to the main benchmark. The
remaining three (FedRecovery, FedEraser, SIFU) were implemented directly
from the published algorithm descriptions.

Table~\ref{tab:baseline-sources} summarises the provenance of each baseline.

\begin{table}[t]
\centering
\scriptsize
\setlength{\tabcolsep}{3pt}
\begin{tabular}{llll}
\toprule
Baseline & Repo & License & Source \\
\midrule
Halimi      & IBM/federated-unlearning  & Apache-2.0 & paper+code \\
FedMemPrune & JWei1999/Rethinking-FU    & MIT        & paper+code \\
FedRecovery & (none)                    & n/a        & paper \\
FedEraser   & Zoybzo/FedEraser-Code     & unlicensed & paper \\
SIFU        & Accenture/Labs-FL (SIFU)  & unlicensed & paper \\
\bottomrule
\end{tabular}
\caption{Baseline sources. ``paper+code'' = algorithm from paper with the central update verified against the authors' code; ``paper'' = no usable public code.}
\label{tab:baseline-sources}
\end{table}

The following hyperparameter choices were made where the published papers
either leave values unspecified or provide values that are not directly
applicable to the shared evaluation protocol.

\paragraph{Halimi.} Ascent steps $=50$ (paper uses ``a small number of
steps'' without a specific value). Projection ball radius $=3$ (paper
reports results across $\{1, 3, 5, 10\}$; $3$ gave the best utility-forgetting
trade-off on the authors' MNIST runs). Local learning rate for ascent
$=0.005$ (paper uses $0.01$ for 10-step runs; proportionally scaled for 50
steps).

\paragraph{FedMemPrune.} Prune quantile $=0.95$ (paper reports a sweep over
$\{0.90, 0.95, 0.99\}$ and recommends $0.95$ for CIFAR-scale benchmarks).
Reset target $=\theta_0$ (paper also considers soft dampening to a scaled
version of the current weight; the hard reset-to-init matches the paper's
default). Pruning applied to all parameters including the classifier head
(paper's default).

\paragraph{FedRecovery.} Weighting $=$ linear in round index (paper
recommends this over uniform, which was also tested and gave worse results).
Subtraction scale $\eta=1$ (paper's default). Sensitivity $C=50$ (matches
the clip norm of the base FedAvg training). Noise scale calibrated to
$(\eps, \delta) = (1.0, 10^{-5})$ via the standard Gaussian mechanism
$\sigma = C\sqrt{2\ln(1.25/\delta)}/\eps \approx 242$.

\paragraph{FedEraser.} Calibration factor $\ell_{\text{calib}} = 0.5$ (paper
recommends between $0.3$ and $0.7$; $0.5$ is the midpoint). History stored
every 5 rounds (paper uses every round; the stride is a storage-cost
compromise forced by the single-M4 evaluation hardware and is noted as a
possible source of degradation).

\paragraph{SIFU.} Target epsilon $\eps^{\star}=1.0$ (paper reports across
$\{0.5, 1.0, 2.0\}$; $1.0$ is the median). Participation rate parameter
$\rho=0.10$ (matches the base training's client sample fraction of 0.20
with a factor-of-two conservatism to account for the influence of a client
that participates and the influence that propagates between rounds).
